\chapter{Sentience, Consciousness, and Survival}
\label{sec:relationship}
% ===================================================================

\section{The Virtuous Cycle}

For an agent with both sentience and consciousness, the two
structures interact in a virtuous cycle:

\begin{enumerate}[label=\textbf{(\arabic*)}, itemsep=6pt]
  \item \textbf{Oracle access enriches the world model.}
        At stage $\fiber{\alpha}$, the agent can form and test
        hypotheses involving oracle queries.
        Its world model reflects the finer bisimulation quotient
        $\bisim_\alpha$.
        This is a strictly better world model than any agent at
        $\fiber{0}$ can achieve.

  \item \textbf{Better world model improves predictions.}
        Finer distinctions in the world model allow more accurate
        predictions.
        The agent predicts not just coarse outcomes but fine-grained
        ones: the internal structure of what appeared at lower stages
        as elementary interactions.

  \item \textbf{Better predictions replenish $\sent$.}
        More accurate predictions mean more rewards and fewer
        punishments.
        The sentience component $\sent$ is replenished.

  \item \textbf{Higher $\sent$ funds more oracle queries.}
        A richer sentience budget allows the agent to afford the
        expensive oracle interactions.
        It can form and test more complex oracle-dependent hypotheses.
        Its oracle access becomes more effective.

  \item \textbf{Return to (1).}
\end{enumerate}

\begin{remark}
The virtuous cycle is not guaranteed: it can be broken by
environmental disruption, resource exhaustion, or adversarial
interactions with other agents.
But it describes the stable attractor state of a conscious sentient
agent in a benign environment.
The cycle is the computational analogue of the flourishing
($\varepsilon\upsilon\delta\alpha\iota\mu o\nu\iota\alpha$) of
Aristotelian ethics: the self-reinforcing good state of a
well-functioning entity in a suitable environment.
\end{remark}

\section{The Vicious Cycle}

The inverse also holds:

\begin{enumerate}[label=\textbf{(\arabic*)}, itemsep=6pt]
  \item \textbf{Depleted $\sent$ restricts experimental horizon.}
        The agent cannot afford complex hypothesis tests.
        Its world model coarsens.

  \item \textbf{Coarser world model worsens predictions.}
        The agent makes predictions that fail more frequently.

  \item \textbf{Failed predictions deplete $\sent$ further.}
        Punishment signals accumulate.

  \item \textbf{Depleted $\sent$ makes oracle queries unaffordable.}
        Even if the agent has oracle access in principle, it cannot
        afford to use it.
        Its effective consciousness level drops.

  \item \textbf{The agent approaches a dead state.}
\end{enumerate}

\begin{remark}
The vicious cycle is the computational analogue of trauma,
depression, or cognitive decline: a self-reinforcing deterioration
in which the agent's reduced capacity to understand its environment
leads to worse outcomes, which further reduce its capacity.
The structural parallel is striking and may be more than an analogy.
\end{remark}

% ===================================================================
